\section{Simulated Business Data Source}
\label{sec:5-1}

Online Retail II is described as a historical extract of data collected from an e-commerce transaction system. The narrative is not "download a dataset and analyze it" but "the analytical system receives transaction records extracted from the sales system".

\section{Simulated E-commerce Transaction System}
\label{sec:5-2}

The simulated architecture contains Order, Product, Customer and Transaction systems. A completed purchase produces invoice, customer, product, quantity, price, timestamp and country information.

\section{Data Extraction}
\label{sec:5-3}

Transaction data are assumed to be extracted from the transactional database into raw analytical storage, followed by validation and analytical transformation.

\section{Collection Frequency}
\label{sec:5-4}

Transactions may be generated continuously while the analytical platform ingests data daily or on another schedule. Ingestion frequency is distinct from model scoring frequency.

\section{Historical Dataset Representation}
\label{sec:5-5}

Online Retail II is the historical extract used for experimentation, and the report transparently identifies it as a public dataset. Properties that affect the design:

\begin{itemize}
    \item It covers roughly \textbf{two years} (December 2009 -- December 2011), which limits the number of labelable weekly snapshots (see Sections~\ref{sec:6-14}, \ref{sec:7-9}).
    \item The retailer is UK-based and many customers are \textbf{wholesalers/businesses}, so spend per customer is highly heterogeneous.
    \item Demand is strongly \textbf{seasonal}, with a peak before Christmas.
\end{itemize}
